\section{Considerações finais}
\label{sec:consideracoes}

A pergunta central foi respondida pela integração de duas camadas parcialmente sobrepostas. O mapeamento bibliométrico identificou evolução temática de eficiência e desempenho para tecnologias digitais, carbono, ciclo de vida e manutenção preditiva; a síntese temática mostrou que as dimensões técnica-operacional, institucional e ambiental e os critérios de desempenho, informação, custo, vida útil, energia e risco formam o núcleo documental, com 109 dos 121 registros temáticos pertencendo também ao estrato bibliométrico de 372. Essa integração fundamentou a matriz de critérios e indicadores e o fluxo de validação para futura parametrização multicritério, respondendo à RQ0 ao ligar evidência, critérios e decisão pública rastreável.

Os resultados sugerem que a lacuna não está na escassez de dados, estruturas ou modelos, mas na cadeia que os conecta a uma decisão institucional validada. Estruturas genéricas, BIM e apoio à decisão foram mais frequentes que métodos multicritério nomeados, e a busca de sensibilidade, ao elevar o aprendizado de máquina de nove para 26 registros, revelou um lote de 17 estudos concentrado em previsão, diagnóstico e integração tecnológica, sem demonstrar uma cadeia decisória completa entre modelo treinado, validação externa, ponderação, vetos e decisão pública auditável. Universidades, \textit{campus} e edifícios públicos foram menos frequentes que edificações genéricas e portfólios, e doze registros apresentaram lacunas específicas para instituições públicas de ensino superior, reforçando a necessidade de validação contextual brasileira.

Diante dessa lacuna, a especificação operacional candidata acrescenta uma matriz rastreável que liga critérios sustentados pela revisão a indicadores, fontes, periodicidades e direções propostos pelo autor, além de um fluxo que ordena validação de conteúdo, normalização, ponderação, agregação e vetos candidatos. A revisão sustenta os critérios; os indicadores permanecem proposições autorais para validação, e a definição de pesos, agregação e limiares dos vetos de risco estrutural, risco à vida, risco de incêndio e descumprimento legal integra a etapa institucional seguinte, dependente de norma, laudo ou autoridade competente.

Essas dependências refletem também os limites do desenho adotado. A cobertura em Scopus, Web of Science e Crossref, entre 2010 e 2026 com o último ano parcial, priorizou metadados bibliográficos interoperáveis e deixou para estudos futuros pré-registro, rastreamento de citações e busca estruturada de literatura cinzenta; triagem e codificação foram conduzidas por um único avaliador, com decisões e regras preservadas para auditoria, e o enriquecimento deliberado da busca de sensibilidade para IA e aprendizado de máquina mede ganho de cobertura temática, não prevalência independente dessas técnicas. O núcleo de 121 registros foi sintetizado predominantemente por título, resumo, palavras-chave e campos estruturados, qualificado por 37 leituras integrais, das quais 19 acrescentaram evidências específicas. A unidade quantitativa é o registro consolidado, e publicações relacionadas podem compartilhar dados ou casos; a classificação do desenho permaneceu indeterminada em 23,1\% dos registros, e a revisão priorizou caracterização documental em vez de avaliação formal de risco de viés.

A continuidade do trabalho deverá começar pela validação de conteúdo, com a participação das áreas de manutenção, engenharia, sustentabilidade, planejamento e orçamento, além dos usuários, na revisão das dimensões, dos indicadores e dos vetos antes da elicitação dos pesos. Em seguida, ordens de serviço, inspeções, custos, consumo, riscos e reclamações deverão receber dicionário, periodicidade e responsáveis, para que o protótipo parametrize e teste o método de elicitação, a normalização, os pesos, a agregação, os limiares e a análise de sensibilidade, comparando estabilidade e utilidade decisória entre unidades e outros portfólios públicos.

Diante da escassez de recursos, esta especificação oferece um caminho possível para justificar cada investimento por critérios documentados e auditáveis, contribuindo para deslocar a priorização da manutenção de uma resposta improvisada para uma prática mais próxima da governança e do valor público, uma vez validada institucionalmente.

\subsection{Disponibilidade de dados}
\label{sec:disponibilidadedados}

Os dados brutos, intermediários e finais que sustentam esta revisão (protocolo de busca, matrizes de triagem e deduplicação, síntese temática e listas de referências em cada estágio) estão depositados em repositório aberto \parencite{oliveira_dadosuplementares_2026}.
